\section{The original local, theta, and operator dictionary}
\label{sec:local-theta-dictionary}

\status{LTD1--LTD8 reconstruct exact transports on the selected objects and
the finite Hecke argument producing the full local rank coverage.  In
particular, an actual \texttt{NBuilt} certificate implies the literal
\texttt{LocsCoverTw}, and eventual such certificates imply
\texttt{LinkEta}.  No additional published Hecke theorem is needed for
this implication: its precise hypotheses occur in \texttt{NdSpec.std}.
The existence of the compatible original geometric, theta, spectral,
and local certificates remains a producer obligation.}

The same field names occur in several different records in this part of
the source.  We distinguish the rational row operators, acting on
\(S.W_L(j)\), from the ambient operators, acting on
\(D.\mathrm{MtL}(\mu_w,j)\).  The former have Hecke indices through
\(r\), and the latter through \(n_V\).  Their comparison must use
the prescribed maps and normalizing scalars.  Merely having the same
name \texttt{hL} does not identify them.

All assertions below concern the pinned original source and the exact
data in their premises.  ``B'' classifies the complete human proof of
the stated implication; it does not classify the unconstructed premise
as elementary.  No new Lean replay is claimed.

% BEGIN LTD1
\paragraph{LTD1: all global selections come from one literal bundle.}
\label{claim:LTD1}
Let \(\mathcal A\) be a family class and let \(\sigma\) be a
\texttt{RawSel Bundle} on that class.  Thus, whenever
\(\mathcal F\in\mathcal A\), \(x>\mathcal F.x_1\), and
\(i\in\mathcal F.I_x\), there is a specified bundle
\[
 b=\sigma(\mathcal F,x,i)
  =(R,S,P,\mathrm{pin},DH,C,B)
\]
on the original \(D=\mathcal F.D(x,i)\), its original object \(O\),
and the same size constant \(c=\mathcal F.C_{\rm sz}\).
Here \(R\) is the tied record, \(S\) its spectral realization,
\(P\) its row family, \(C\) the compatible columns, and
\[
 B:\texttt{RealTB}
   (R,P,C.\mathrm{face},C.\mathrm{Sp},
                 \mathrm{zsc}(P,DH,c),S).
 \tag{LTD.1}
\]
The source's selections are exactly
\[
 \begin{split}
 \mathrm{sel}(\mathcal F)&=
   \bigl\langle\mathcal F.x_1,
           (x,i)\longmapsto b.R\bigr\rangle,\\
 \mathrm{selS}&=b.S,\quad \mathrm{selP}=(b.P,b.\mathrm{pin}),\quad
 \mathrm{selDH}=b.DH,\quad \mathrm{selC}=b.C,\quad
 \mathrm{selB}=b.B.
 \end{split}                                                    \tag{LTD.2}
\]
Indeed, substituting these projections in the definition of
\texttt{ZscAt} gives the very scale in (LTD.1): its last component
is set to zero and its previous component is changed by
\(-c\,z_{6,u}\).  All other terms use the same \(P,DH,D,O.L\).
Consequently the dependent type of \(b.B\) is precisely the type
required by \texttt{SelB}; no new compatible realization is chosen.

In particular, the indefinite realization selected by
\texttt{selIOf} is exactly \(b.B.Q\), and its coefficient field
selected by \texttt{selKOf} is exactly \(b.B.T.Q.c.F\).
The source's \texttt{wire\_sel} uses this one \(Q=b.B.Q\) for its
all-tuples comparison, its point comparison and its base-rank
comparison.  Likewise \texttt{link\_alt} takes its spectral
instance and its instance tie from \(b.S\), its real volume tie
from \(b.R.T.\mathrm{theta.real}\), and its scale tie from the
original \(\mathcal F\).  These are common-object identities, not
assertions inferred from equal dimensions.

The condition \(x>\mathcal F.x_1\) already implies the weak
threshold requirement \(\mathcal F.x_1\le x\) in (LTD.2).
No further eventual threshold is needed to align these projections.
This proves the compatibility of the selection once \(\sigma\)
is supplied; it does not construct the bundles.
% END LTD1

% BEGIN LTD2
\paragraph{LTD2: the local installation preserves the original theta and form objects.}
\label{claim:LTD2}
Fix \(D\) and \(L:\texttt{LocAll}(\texttt{lay0}(D))\), and put
\(D_L=\texttt{withLoc}(L)\).  The local installation replaces the
local index set, local rings, modules, projections, automorphic
specializations and associated local comparisons by the fields of
\(L\).  It also sets \(Z,Z_{\rm good}\) to the specified actual
subsets and replaces the row number \(\eta_J\), for \(1\le j\le m\),
by
\[
       \eta_{J,L}(j,w)=\sum_{a\in\mathrm{locs}_L(j,w)}\eta_{\rm loc}(a).
       \tag{LTD.3}
\]
Outside that range it preserves the previous \(\eta_J\).
It preserves the host, its coefficient fields and their embeddings,
all row spaces \(\mathrm{MtL},\mathrm{Mt},\mathrm{Mj}\), their
original realization maps and lattices, and their theta maps.

Here is a direct check of the claim.  \texttt{LocNd.nd} changes
exactly the local-indexed fields of \texttt{NData}.  In
\texttt{LocTw.ntw}, only the maps from these new local indices to
the old twist indices, their local coefficient maps, and the local
characters are replaced.  The twist indices, their rings \(O_0,O,O_k\),
the two side embeddings, the universal Artin maps, the twist
characters, signatures and toric theta values are copied.  In
\texttt{LocAll.lay}, \texttt{nexp} and \texttt{nth} are record
copies.  The subsequent \texttt{fixZ} changes only the two finite
sets; \texttt{reEta} changes only \(\eta_J\) in the underlying
row datum.  Therefore every claimed carrier and every claimed map
has the same defining expression after the necessary dependent
type identifications.  In particular, the coefficient-valued
\(\mathrm{th},\mathrm{thPsi},\mathrm{thU},\mathrm{thP}\), the
form action and Hecke operators, and the actual evaluation maps
of \texttt{NTh} are the originals.

For an original form family \(M\), normalized choice \(C\),
character frame \(cf\), and tuple tie \(tp\), the source defines
\[
 M_L=\texttt{mtL}(L,M),\qquad C_L=\texttt{chL}(L,C),\qquad
 cf_L=\texttt{cfL}(L,cf),\qquad tp_L=\texttt{tpL}(L,tp).
 \tag{LTD.4}
\]
These copy the old fields, including \(C.\mathrm{twI},C.\mathrm{sd},
C.Q,C.k',C.\psi',C.\mathrm{embW},C.e_{Mt}\), and the tuple
characters.  Thus the exact semilinear basis and lattice equalities
in \texttt{NChoice} retain their maps, not merely their numerical
ranks.  This is the source's \texttt{Kept} comparison at this step.

Preservation is an implication on an existing object.  It does not
turn an initialized zero theta map into a nonzero one.  In particular,
the separate \texttt{NTh.ThEv} and \texttt{NTh.EmbBRange}
propositions are not fields proved merely by presenting an
\texttt{NTh} record.  Whenever such a proposition has been proved
for the original maps, the equality just checked transports it
unchanged; its mathematical content has not been manufactured by
the installation.
% END LTD2

% BEGIN LTD3
\paragraph{LTD3: integral simultaneous generalized eigenspaces.}
\label{claim:LTD3}
Let \(O\) be a discrete valuation ring, \(K\) its fraction field,
and \(V\) a finite-dimensional \(K\)-vector space.  Let
\(\mathcal S\subseteq\operatorname{End}_K(V)\) be a possibly
infinite pairwise commuting family.  Suppose each member splits
over \(K\), in the sense that its generalized eigenspaces span
\(V\), and every eigenvalue having a nonzero generalized
eigenspace lies in \(O\).  Let \(T=O[\mathcal S]\).  For an
\(O\)-algebra homomorphism \(\psi:T\to O\), define
\[
 W_\psi=\bigcap_{t\in T}\operatorname{GEig}(t,\psi(t)),
 \qquad
 \mathfrak m_\psi=\ker(T\xrightarrow{\psi}O\longrightarrow O/\mathfrak m_O).
 \tag{LTD.5}
\]
The scalar inclusion \(O\hookrightarrow K\) is implicit in
\(\operatorname{GEig}\).  There are only finitely many nonzero
\(W_\psi\), and
\[
                     V=\bigoplus_{W_\psi\ne0}W_\psi.             \tag{LTD.6}
\]

We prove this without a finiteness assumption on \(\mathcal S\).
Begin with \(V\).  If a nonzero current block is split into two
or more generalized eigenspaces by some \(s\in\mathcal S\),
perform that refinement.  Every new block is invariant under
all of \(\mathcal S\), since commuting operators preserve the
kernels of powers of \(s-\lambda\).  The restriction of another
split operator to an invariant subspace still splits: its
generalized spectral projections are polynomials in that
operator, hence preserve the invariant subspace.  Each strict
refinement increases the number of nonzero summands, which is
at most \(\dim_KV\).  We therefore reach a finite direct sum
\(V=\bigoplus_aV_a\) such that every \(s\in\mathcal S\) has on
each nonzero \(V_a\) exactly one generalized eigenvalue
\(\lambda_a(s)\in O\).

On \(V_a\), write \(s=\lambda_a(s)I+N_s\).  Each \(N_s\)
is nilpotent, and all these nilpotents commute.  A polynomial
with zero constant term in finitely many commuting nilpotents
is nilpotent: in a sufficiently high power, every monomial has
an exponent at least the nilpotence exponent of one factor.
Consequently every polynomial expression \(t\in T\) acts as
\(\lambda_a(t)I+N_t\), where \(N_t\) is nilpotent and the scalar
is the same polynomial in the \(\lambda_a(s)\) and its
\(O\)-coefficients.  This scalar is independent of the expression
for \(t\).  Indeed, a relation would make \(cI\) nilpotent for
the difference \(c\) of two putative scalars; on the nonzero
space \(V_a\), that implies \(c=0\).  Thus \(\lambda_a\) is an
\(O\)-algebra homomorphism \(T\to O\).

Two distinct blocks were separated at a refinement and hence
have distinct values on some generator.  Their characters
\(\lambda_a\) are distinct.  Their generalized character spaces
are exactly the blocks: on any other block, a generator with
a different value has invertible scalar-plus-nilpotent difference.
The same observation proves that any nonzero \(W_\psi\) is one
of these blocks.  This proves (LTD.6).  For \(V=0\), all spaces
in (LTD.5) are zero and the conclusion is immediate.

For a maximal ideal \(\mathfrak m\subset T\), put
\[
 V_{\mathfrak m}=\sum_{\mathfrak m_\psi=\mathfrak m}W_\psi,
 \qquad V^{\rm off}_{\mathfrak m}
            =\sum_{\mathfrak m_\psi\ne\mathfrak m}W_\psi.        \tag{LTD.7}
\]
Each \(\mathfrak m_\psi\) is maximal: residue evaluation is
surjective because it contains the images of all constant
elements of \(O\).  Therefore, for any finite set \(\mathcal M\)
containing every maximal ideal of \(T\), grouping (LTD.6) by
these residue kernels gives
\[
 \sum_{\mathfrak m\in\mathcal M}\dim_KV_{\mathfrak m}=\dim_KV,
 \qquad
 \bigcap_{\mathfrak m\in\mathcal M}V^{\rm off}_{\mathfrak m}=0.
 \tag{LTD.8}
\]
Ideals with no supported character contribute a zero summand.
The argument does not assert diagonalizability, and does not
need to prove that \(T\) is finite as an \(O\)-module.
% END LTD3

% BEGIN LTD4
\paragraph{LTD4: the actual local specialization gives complete rank coverage.}
\label{claim:LTD4}
Retain \(D,L,M,C\) from LTD2 and a specified
\[
 \mathrm{Sp}:\texttt{NdSpec}(\mathrm{Pl},M,C,\mathrm{big},L.X,L.T).
 \tag{LTD.9}
\]
For every \(1\le j\le m\) and every original place \(w\) with
\(D.\mu_{\rm At}(w)\ne D.\mu_0\), one has
\[
 \sum_{a\in D_L.\mathrm{locs}(j,w)}
          \operatorname{rank}_{O_a}M_a
       =\dim_{D.L}D.\mathrm{MtL}(D.\mu_{\rm Tw}(w),j).
 \tag{LTD.10}
\]
Moreover the original projections \(\mathrm{pr}_{Mt,a}\), after
the explicit index casts, are jointly injective on the right-hand
space.  Equation (LTD.10), for these exact quantifiers, is
\texttt{LocsCoverTw} on \(D_L\).

Fix \(j,w\).  Take \(O=\mathrm{Sp}.B.O_c(j,w)\), its fraction
field \(K\), \(V=\mathrm{Sp}.B.V(j,w)\), and the family
\(H_{v,k}=\mathrm{Sp}.B.h_K(j,w,v,k)\) with the literal
\texttt{GenAt} restrictions: \(v\) is split, away from the
coefficient prime, the specified auxiliary set and \(Q_c(j,w)\),
and \(1\le k\le n_V\).  The field \texttt{Sp.std} gives
pairwise commutation and integrality of every occurring
generalized eigenvalue.  The field \texttt{Sp.B.split} gives
the splitting hypothesis of LTD3.

The original row-basis axiom first makes
\(D.\mathrm{MtL}(\mu_w,j)\) finite dimensional: its basis maps
to a linearly independent family in the finite-dimensional
original complex row space.  The field \texttt{iV\_basis} maps
that same basis to a \(K\)-basis of \(V\).  Thus
\[
 \dim_KV=\dim_{D.L}D.\mathrm{MtL}(\mu_w,j),\qquad
                 i_V\text{ is injective}.                         \tag{LTD.11}
\]
For the second assertion, write a vector in the original basis.
A zero image gives a zero combination in the image basis, so
each coefficient has zero image under the field embedding and
is zero.

The fields \texttt{mx\_max}, \texttt{mx\_inj}, and
\texttt{mx\_all} identify the finite set
\(D_L.\mathrm{locs}(j,w)\) bijectively with the maximal ideals
used in (LTD.8).  For each index \(a\), the specified map
\[
 \Psi_a:K_a\otimes_{O_a}M_a\longrightarrow V
\]
is semilinear for the field isomorphism \(\mathrm{phK}_a:K_a\simeq K\),
is injective, and has image exactly \(V_{\mathrm{mx}(a)}\).
Transporting a basis along this isomorphism gives
\[
 \operatorname{rank}_{O_a}M_a
       =\dim_KV_{\mathrm{mx}(a)}.                                 \tag{LTD.12}
\]
Here \(M_a\) is finite free over the given DVR, so extending a
free basis to its fraction field preserves its length.
Summing (LTD.12), using injectivity and exhaustiveness of
\(\mathrm{mx}\), then (LTD.8) and (LTD.11), proves (LTD.10).

Finally suppose every original local projection of \(y\) is
zero.  Membership \(a\in\mathrm{locs}(j,w)\) supplies the
identities of the row number, place and twist index; these give
the source casts \texttt{castMtInv}.  The actual
\texttt{Psi\_pr} equality says
\[
 \Psi_a(\mathrm{pr}_{Mt,a}(y))-i_V(y)
                    \in V^{\rm off}_{\mathrm{mx}(a)}.
 \tag{LTD.13}
\]
Its first term is zero.  By (LTD.8), \(i_V(y)=0\), and (LTD.11)
gives \(y=0\).  This proves joint injectivity on the original
space, not on an arbitrarily chosen space of the same dimension.

No \texttt{AutSpec} hypothesis is needed for this proof.
If one wants the entire \texttt{NormalizedLocal}\allowbreak\texttt{Compatibility}
record, \texttt{Sp} supplies all its displayed algebraic fields;
the two additional equalities \texttt{AutSpec.iotaQ} and
\texttt{iotaQC} turn \texttt{Sp.iotaC\_coh} into the required
complex embedding square.  This is exactly the \texttt{cover}
construction in \texttt{heckeOpen\_withLoc\_alt}.
% END LTD4

% BEGIN LTD5
\paragraph{LTD5: the exact original certificate implies the terminal coverage link.}
\label{claim:LTD5}
For every original datum \(D'\), every fixed local Langlands
input \(h_{LL}\), and every supplied
\(N:\texttt{NBuilt}(h_{LL},D')\), one has
\[
                    \texttt{LocsCoverTw}(D').                    \tag{LTD.14}
\]
Indeed, \(N\) contains \(D,L,M,C,\mathrm{Sp}\) as in LTD4 and
the equality \(D'=\texttt{withLoc}(L)\).  Substitute this
equality and apply LTD4.  The proof never applies the theorem
denoted by \(h_{LL}\); it uses the already supplied local
certificate.  The role of local Langlands in producing that
certificate is separate.

The stronger \texttt{NTrue} record additionally identifies the
entire dependent tuple
\[
 (D',cf',C',tp')
       =(D_L,cf_L,C_L,tp_L).                                     \tag{LTD.15}
\]
Thus (LTD.14) is compatible with the exact characters and choice
used later.  In \texttt{RowsAtStart.}\allowbreak\texttt{nTrueAtWeight}, both the datum
equality and (LTD.15) are reflexive: the constructor packages
these very objects.  No isomorphic replacement is inserted.

Consequently, for a fixed \(\mathcal F:\texttt{FamP}(\epsilon)\),
the sufficient family condition is precisely
\[
 \exists x_0\ \forall x\ge x_0\ \forall i\in\mathcal F.I_x,
       \exists N:\texttt{NBuilt}(h_{LL},\mathcal F.D(x,i)).
 \tag{LTD.16}
\]
Applying (LTD.14) separately at each \((x,i)\), without changing
\(x_0\), gives the source's \(\texttt{LinkEta}(\mathcal F)\).
When the producer requires \(x>\mathcal F.x_1\), replace
\(x_0\) by \(\max(x_0,\mathcal F.x_1+1)\) before applying it.
If (LTD.16) is provided for every member of the actual selected
class, the same assertion is \texttt{Closed} on that class,
and on its \(D3\) subclass by the defining inclusion.

This is a strict reduction of the coverage consumer to the
original \texttt{NBuilt}/\texttt{NdSpec} producer.  It does not
assert that an arbitrary \texttt{TSel} carries an \texttt{NBuilt}
certificate.  The generic source theorem \texttt{link\_eta}
uses only the selected tied records' \texttt{MtFam},
\texttt{NormalizedLocalCompatibility} and \texttt{HeckeTied};
its \texttt{Duty} argument is unused.  That generic route calls
\texttt{heckeStd}.  The concrete route (LTD.16) instead retains
the \texttt{std} proof already present in the actual local
construction, and therefore introduces no second independent
Hecke integrality supplier.
% END LTD5

% BEGIN LTD6
\paragraph{LTD6: the original row spectrum and its normalized Hecke action.}
\label{claim:LTD6}
Fix one source geometric state \(x:\texttt{St}(P)\) and a
compatible \(l:\texttt{LC}(h_{LL},x)\).  This is the source's
\texttt{StrictHenselian}\allowbreak\texttt{AcyclicLevel} in the expanded naming.
Its rational form datum is
\[
 S=\texttt{ratDatR}(x.m.z.R,x.m.z.cG,l.\mathrm{rat}.k_{K_s},
                                                l.\mathrm{rat}.hk).
\]
Under the installation \(L\), the source chooses
\[
 D_L=\texttt{DL}(x,L)=\texttt{withLoc}(L),\qquad
 \texttt{datL}=\texttt{Dat.transfer}(x.m.z),
\]
using the identity finite-group comparison.  Consequently
\[
 S_L=\texttt{ratL}(L,S),\quad H^{\rm row}_{L,j,k,v}
              =l.\mathrm{rat}.\mathrm{In}.h_L(j,k,v),\quad
 e_{L,j,\tau,k,v}=l.\mathrm{rat}.\mathrm{In}.\mathrm{eig}(j,\tau,k,v).
 \tag{LTD.17}
\]
These are equalities of the operators on the transported original
spaces.  Their proof is substitution in \texttt{SL},
\texttt{ysL}, \texttt{hkH}, \texttt{hkHL}, and \texttt{eigD}:
\texttt{SL\_heq} and \texttt{hL\_eq} reduce to reflexive
equalities.  A final row reading \(Y\) is related by the specified
\texttt{YRel}, which keeps its actual representation and its
\(e_W,j_{Mt}\) readings.  \texttt{normalizedRowDatum} uses that
relation to identify its representation with
\(\tau\mapsto\texttt{rowRepR}(R.\theta,Y.e_{FT}(\tau))\).

To check that this identification is useful for the spectral
consumer, write \(G_j\) for \texttt{goodL}: the intersection of
the eigenspace of the original finite-group projector with the
specified generic-character eigenspaces.  All operators and
characters defining this intersection are preserved in LTD2 and
(LTD.17), so its transported subspace is literally the same
\(G_j\).  The original \texttt{RatIn.span} assertion is
\[
 G_j\ \subseteq\sum_{\substack{\tau:\,\mathrm{isot}(j,\tau)\ne0\\
                                      \mathrm{goodRow}(\tau)}}
 \left(G_j\cap\bigcap_{\substack{1\le k\le r\\
                      v\text{ split},\ v\notin S_b}}
       \ker(H^{\rm row}_{j,k,v}-e_{j,\tau,k,v})\right).
 \tag{LTD.18}
\]
Substitution proves exactly (LTD.18) for the final original row
model.  Likewise, if \(\Lambda_j\) is the
\(\mathcal O_{D.L}\)-span of its prescribed integral coefficient
module, the original \texttt{hL\_int} becomes, with the same
quantifiers and the same \(S_b\),
\[
 \forall j,k,v\ \exists n\quad
     (\operatorname{N}v)^n H^{\rm row}_{j,k,v}
                    (\Lambda_j\cap G_j)\subseteq\Lambda_j,
 \tag{LTD.19}
\]
for \(1\le j\le m\), split \(v\notin S_b\), and \(k\le r\).
No uniform denominator exponent is inferred: \(n\) may depend
on \(j,k,v\), exactly as in the original field.

The coefficient-valued Hecke action is related to this row
action by the actual map \(e_F\) and actual unit \(u_{i,b,v}\):
\[
 e_F H^{\rm exp}_{i,j,b,k,v}
       =u_{i,b,v}^{\,k}(1\otimes H^{\rm row}_{j,k,v})e_F,
       \qquad 0\le k\le r,\quad v\notin S_0(i),\ v\text{ split}.
 \tag{LTD.20}
\]
This follows from \texttt{SideB.hkEq} and the equality
\(l.b.YS.e_F=l.b.Bs.e_F\); \texttt{hkH} defines the scalar as
the \(k\)-th power of the embedded \texttt{cHk} unit.  For
\(k=0\) it is one.  For any scalar \(a\), induction on \(n\)
in (LTD.20) gives
\[
 e_F(H^{\rm exp}-u^ka)^n
       =u^{kn}(1\otimes H^{\rm row}-a)^n e_F.                    \tag{LTD.21}
\]
Thus \(e_F\) carries both ordinary and generalized eigenspaces
to those with the prescribed rescaled eigenvalue.  This proves
the exact spectral transport; it does not assert the initial
spanning or integrality in (LTD.18)--(LTD.19) without
\texttt{RatIn}.

There is a second family \(\mathrm{Sp}.B.h_L/h_K\), used in
LTD4, on \(\mathrm{MtL}(\mu_w,j)\) with \(k\le n_V\).
Its operator square is \texttt{NLocBase.hL\_def/hK\_def}, and
its local action square is \texttt{NdSpec.act\_hK}.  These
are not (LTD.17).  In particular,
\texttt{ChoiceSide.eMt\_level} only gives an extending complex
embedding, equality of the level subgroup, and equality of its
character after \texttt{kMapE}.  It contains no pointwise
equality for \(e_{Mt}\) and no Hecke square.  The missing
comparison cannot be supplied by silently identifying these
two operator families.
% END LTD6

% BEGIN LTD7
\paragraph{LTD7: localization transports the existing theta evaluation and export series.}
\label{claim:LTD7}
Continue with the same \(x,l,L\) as in LTD6.  The source's
\texttt{ModuliPointsState} contained in \(l\) supplies an actual
\texttt{GeoTie.ThetaTie} between its original expansion and
\(D.\mathrm{nth}\).  Copy its points, covering maps, evaluation
identities, unramified forms, Hecke identities and unramified
extension condition.  By LTD2 the coefficient fields,
evaluation maps and Hecke operators appearing on both sides
are unchanged, and by LTD6 their real model has been transported
by the identity comparison.  Every copied equality therefore
has exactly the required target and source.  This constructs
the final \texttt{ThetaTie}; it is the source's
\texttt{geoL}.  In particular, it uses the existing evaluation
identity, rather than deducing one from the existence of a
theta-named record.

There is also an exact transport for the exported series.  Let
\(X'\) be any \texttt{NExpTied} on \(D_L\), and suppose
\[
      h:\texttt{ExpRel}(L,l.b.Bs.\mathrm{exp},X').                 \tag{LTD.22}
\]
This comparison has precisely three fields.  It transports
\texttt{TauTrue}; it preserves the ordered pair of bad-place
sets \((QP,DP)\) of every admissible level; and it transports
every \texttt{Reads} factorization.  The latter says that a
specified pointwise reading \(r(f,g)\) factors through the
actual realization by an injective semilinear coefficient map:
\[
 r(f,g)=\omega((\mathrm{real}(f))(g)),\qquad\omega\text{ injective}.
 \tag{LTD.23}
\]
No equality of the two realizations themselves is presumed.

Construct the localized export by retaining the original
characters, finite class sets and representatives, measures,
Hermitian frames, integral lattices, additive character and
normalizing scalars.  Only their ambient datum is relabeled.
The twist maps in the two directions are inverse identity
readings of \((Q,\chi,\mathrm{res})\); thus all original
admissibility and finite-class conditions have the same
meaning.  The \texttt{setupL} construction has the following
pointwise equalities under this relabeling:
\[
 \theta_L=\theta,\qquad \mathrm{gens}_L=\mathrm{gens},\qquad
 \xi_L=\xi,\qquad \mathrm{classSum}_L=\mathrm{classSum},
 \tag{LTD.24}
\]
and its predicates \texttt{IsDatum} and \texttt{IsHeckeDatum}
are equivalent to the originals.  To see this, expand a theta
or class-sum vector into its original frame, place map, finite
representatives and vector function.  Each of these has just
been retained.  The datum predicates assert splitting,
avoidance, finiteness, support and membership for those same
arguments; the Hecke predicate asserts the same chosen element,
representatives and determinant.  Each clause is unchanged.

An original \texttt{CellTie} can now be copied field by field.
Its support, class, polynomial coefficient and \(\mathrm{rdV}\)
maps are kept.  Its datum-surjectivity and theta equalities use
the first equality in (LTD.24), its twisted-generator equality
uses the second, its universal coefficient identities and
off-support vanishing use the third, and its theta-projection
identity uses the fourth.  The equality of integral lattices
and injectivity of \(\mathrm{rdV}\) refer to the unchanged maps.
The finite-group and level-character equations are unchanged
as well.  This checks every field of \texttt{CellTie},
including its conditional antecedents.

Keep exactly the original normalization scalar \(\nu\ne0\).
The norm assertion \texttt{IsNorm} transports in both directions
because its actual family, measures and maps have been retained.
The three rationality conditions on the base theta values,
twisted generators and universal coefficients are equivalent
by (LTD.24).  Hence both conditional clauses in
\texttt{IsScalar}, as well as \(\nu\ne0\), transport.  This
preserves the original scalar choice even when a conditional
rationality antecedent is false; it does not add an unconditional
nonvanishing theorem.  The copied Hecke tie follows from the
\texttt{IsHeckeDatum} equivalence, and the universal-character
tie is retained by LTD2.  We have constructed the localized
\texttt{ExportTie} with its original normalization.

To finish \texttt{SerAt}, retain its bad-place equality, its
selected-place equality, the bounded-lattice equivalence, the
off-range zero condition and its original lambda datum and
lambda ties.  Its only remaining field, \texttt{join}, is the
factorization (LTD.23) for the pointwise expression
\[
                 \mathrm{rdV}(f)(\mathrm{finHom}(\mathrm{PV}(g))).
\]
The old \texttt{join} supplies that reading; apply
\(h.\mathrm{reads}\) to it.  This produces its required
injective semilinear map for \(X'\), on the same \(f,g\).
We have therefore proved
\[
 \texttt{ExpRel}(L,l.b.Bs.\mathrm{exp},X')
  \Longrightarrow
 \operatorname{Nonempty}\bigl(\texttt{SerAt}
     (\mathrm{ysL},\mathrm{datL}.\mathrm{geo}.v_0,
                     X',\mathrm{pvL},\mathrm{cfD},\mathrm{pqD})\bigr).
 \tag{LTD.25}
\]
The original \texttt{ThetaTie}, \texttt{SerAt} and the comparison
(LTD.22) remain explicit premises.  Their analytic rationality,
convergence and initial realization content is not replaced
by the transport proof.
% END LTD7

% BEGIN LTD8
\paragraph{LTD8: the coefficient field is the actual embedded field.}
\label{claim:LTD8}
Fix the source's real choice \(ch\).  Its predicate
\texttt{FieldOK} means \texttt{NumbersIn} for the particular
embedding \(\iota=D.\mathrm{embLL}:D.L\hookrightarrow\mathbb C\).
The required numbers comprise the specified character values,
eighth roots, the place-dependent roots of unity, phases,
square roots of place norms, finite expansion jets, critical
toric ratios and normalized bottom sums.  Thus, for any field
embedding \(\sigma:K\hookrightarrow\mathbb C\),
\[
 \texttt{FieldIn}_\sigma(ch)\quad\hbox{and}\quad
 \sigma(K)\subseteq\iota(D.L)
       \quad\Longrightarrow\quad\texttt{FieldOK}(ch).            \tag{LTD.26}
\]
For every required number, compose its asserted membership in
\(\sigma(K)\) with the displayed containment.  In the bottom
sum clause retain the same triangles and the same nonzero
period scalars, and apply the containment to each normalized
sum.  All quantifiers over tuples, supported expansions and
places are untouched.  This proves every clause of
\texttt{NumbersIn}, hence (LTD.26).

Conversely, \texttt{FieldOK} implies \texttt{HasNumberField}:
take the image field \(\iota(D.L)\subseteq\mathbb C\).
The embedding identifies it with the original number field,
so it has the same finite dimension over \(\mathbb Q\);
every required number is in that image.  The converse needs
the containment in (LTD.26); existence of some number field
alone gives no such containment for the fixed datum.

The jets clause immediately makes the specified top and row
jets rational at their actual pairs of points, and the
character clause gives \texttt{ValLL}.  For the further
\texttt{FieldCond} consumer the source also requires the same
real additive character and theta datum, equality of the two
grid characters at every relevant index, equality of the
selected place set and its depths, nonvanishing of the
denominator Hecke values at the specified positive integers,
existing bottom triangles, and integrability of their two
kernel expressions.  Under those additional premises,
substitution gives the required roots and character values;
the critical-ratio clause supplies
\(\iota(\rho)=\Lambda(s)/\Lambda_0(s)\), and the nonzero
denominator gives \(\Lambda(s)=\iota(\rho)\Lambda_0(s)\).

For the last clause, the existing bottom-triangle witness
activates \texttt{BottomIn}.  This supplies a possibly newly
chosen triangle \(X\) satisfying the same \texttt{IsBottomTri}
conditions, and nonzero complex periods \(\Omega,\Omega'\).
Its \texttt{SumsIn} certificate already asserts that every
normalized monomial integral \(I/\Omega\), respectively
\(I'/\Omega'\), lies in the actual field image, conditional on
integrability of its corresponding kernel.  The additional
\texttt{hI} hypothesis is quantified over every such triangle,
so it applies to this selected \(X\) and discharges both
integrability guards.  For every allowed \(t\), finite-adelic
point \(g\), and respective monomial indices \(k,k'\), there are
therefore
\(z,z'\in D.L\) with
\[
 \iota(z)=\frac{I(t,k,g)}{\Omega},\qquad
 \iota(z')=\frac{I'(t,k',g)}{\Omega'},\qquad
 I(t,k,g)=\Omega\iota(z),\quad
 I'(t,k',g)=\Omega'\iota(z').
 \tag{LTD.27}
\]
The last two equalities follow by multiplying through the
nonzero periods.  Together with the integrability statements
they are exactly \texttt{SumsIn.rationalSums}; the remaining
triangle fields transport by the stated character, place and
rank equalities.  Rationality of these integrals is thus an
explicit part of the \texttt{SumsIn} input.  Without these
premises, \texttt{FieldOK} is not the whole period formula.
% END LTD8

\subsection*{The reduced producer boundary}

The following distinctions are consequences of the proofs above.
They specify remaining work on the original objects, rather than
introducing replacement models.

\begin{longtable}{p{0.22\textwidth}p{0.34\textwidth}p{0.36\textwidth}}
\textbf{Consumer}&\textbf{Completed implication}&\textbf{Retained producer}\\\hline
Common selected objects&LTD1: every selection is a projection of the
same bundle&The common bundle/compatible design existence theorem.\\
Local installation&LTD2: original forms, theta maps, embeddings,
choices and operators survive the installation&The original
\texttt{LocAll} and the earlier nonplaceholder theta realization.\\
Local rank coverage&LTD3--LTD5: actual \texttt{NdSpec}/\texttt{NBuilt}
gives \texttt{LocsCoverTw}; eventual certificates give
\texttt{LinkEta}&Actual local specialization on the common model,
including \(\Psi\), its range and its projection square.\\
Actual row spectrum&LTD6: the specified rational row operators,
eigenvalues, lattices and normalized coefficient action are transported
exactly&Original \texttt{RatIn.span}, \texttt{hL\_int}, row reading,
and the compatible geometric state.\\
Theta evaluation and series&LTD7: existing \texttt{ThetaTie} and
\texttt{SerAt} survive localization with their original scalar&The
original evaluation/series proofs and \texttt{ExpRel}; the separate
\texttt{EMtTheta} comparison and actual row theta square.\\
Coefficient rationality&LTD8: every inclusion and field transport is
on the specified embedding&Actual toric-ratio/jet/period values,
their prescribed field containment and analytic hypotheses.\\
\end{longtable}

The next reconstruction should construct these retained inputs at
the state used by \texttt{RowsAtStart}.  They must then be retained
through the bundle constructor.  In particular, the local coverage step no
longer warrants a separate deep reconstruction after the actual
\texttt{NdSpec} has been supplied.  The shared local/theta dictionary
has been narrowed by exact transport theorems, while its substantive
existence content remains open in the human reconstruction.
